\section{Conclusiones}

El desarrollo del presente proyecto ha permitido alcanzar de forma satisfactoria
los objetivos planteados inicialmente, logrando no solo una adquisición de conocimientos en 
el ámbito de la creación y edición de vídeos, sino también la capacidad de aplicar
dichos conocimientos en un contexto práctico y realista.


En primer lugar, el análisis exhaustivo y la comparativa de las distintas herramientas
del mercado audiovisual ha revelado las fortalezas y limitaciones de las opciones actuales:
\begin{itemize}
    \item \textbf{Canva Video}: Se presenta como una opción ágil para el diseño rápido 
    apoyada en la filosofía \textit{drag-and-drop}. Ha resultado ser más útil de lo 
    esperado, pero es cierto que resulta insuficiente frente a ediciones de vídeo complejas.
    \item \textbf{DaVinci Resolve}: El análisis ha consolidado a esta herramienta como la
    opción idónea al ofrecer un perfil altamente profesional. Permite realizar edición no 
    lineal, corrección de color y aplicar efectos visuales de forma muy superior a la
    competencia. Sin embargo, su curva de aprendizaje es considerablemente más pronunciada,
    se requiere de un hardware potente para su uso óptimo y la versión gratuita presenta
    más limitaciones de las que se pueden ver a simple vista.
\end{itemize}

En segundo lugar, la aplicación práctica mediante la creación del tráiler de vídeo ha 
puesto a prueba los conocimientos teóricos adquiridos. La grabación, realizada sin depender
de equipos especializados y con un presupuesto nulo, ha demostrado que es posible obtener 
resultados de calidad utilizando recursos accesibles. También ha ayudado a seguir
desarrollando las habilidades organizativas, de planificación y de trabajo en equipo, 
aspectos fundamentales en esta clase de proyectos, pues pueden ser la diferencia entre el
 éxito o el fracaso de un proyecto audiovisual.

En definitiva, la integración de todos los elementos desarrollados —desde el análisis
comparativo de herramientas hasta la ejecución técnica del tráiler— ha proporcionado una 
visión global y multidisciplinar del sector audiovisual. Los resultados obtenidos confirman 
que, con una base teórica sólida y un uso estratégico de los recursos digitales actuales,
es posible transformar una idea creativa en un producto óptimo sin necesidad 
de grandes infraestructuras. Este aprendizaje no solo consolida competencias técnicas en 
edición, sino que nos ha proporcionado una visión crítica sobre la comunicación multimedia que será de
gran valor en el futuro.